% ==========================================================
% 050_forecast_admissibility_surface.tex
% Forecast Admissibility Surface (FAS)
% ==========================================================

\section{The Forecast Admissibility Surface}
\label{sec:fas_surface}

The Forecast Admissibility Surface (FAS) formalizes admissibility as a
deterministic classification over slice domains derived from baseline error
anatomy. Rather than producing a scalar score or ranking, FAS defines a
\emph{decision surface} that partitions the forecast domain into regions where
forecasting and readiness control are permitted, constrained, or disallowed.

Let $\Domain$ denote the forecast domain under consideration and let
$\{\Slice_k\}_{k \in \SliceKeys}$ denote the partition of $\Domain$ induced by the
chosen slice keys. FAS defines a mapping
\[
\FAS : \Slice_k \;\longrightarrow\; \{\Allowed, \Conditional, \Blocked\},
\]
where each slice is assigned exactly one admissibility class based on observable
baseline error anatomy.

\subsection{Partition of the Forecast Domain}

The admissibility surface induces a three-way partition of the domain,
\[
\Domain = A \;\cup\; C \;\cup\; B,
\]
where $A$ denotes the set of \Allowed\ slices, $C$ the set of
\Conditional\ slices, and $B$ the set of \Blocked\ slices. This partition is
exhaustive and mutually exclusive: every slice belongs to exactly one region.

Crucially, FAS does not impose a total ordering or notion of “degree” of
admissibility. The surface is categorical by design. A slice is either admissible
without restriction, admissible only under constrained conditions, or
inadmissible for forecasting and control.

\subsection{Deterministic Rule Closure}

Admissibility classification is determined by a fixed set of thresholded rules
applied to baseline error anatomy. These rules operate on support size, spike
frequency, and tail magnitude, and are evaluated independently for each slice.
Given a set of thresholds, the mapping from error anatomy to admissibility class
is fully deterministic.

This determinism is a central design principle. For a fixed dataset, baseline
model, slice definition, and threshold configuration, FAS will always produce
the same admissibility surface. No learned parameters, stochastic procedures, or
optimization routines are involved. As a result, admissibility decisions are
fully reproducible and auditable.

\subsection{Surface Interpretation}

The term \emph{surface} is used deliberately. FAS should be interpreted as a mask
laid over the forecast domain that constrains downstream behavior. It does not
modify forecasts directly, nor does it intervene in model training. Instead, it
controls \emph{where} such activities are permitted to occur.

Slices in the \Allowed\ region may flow freely through modeling, evaluation, and
readiness-control pipelines. Slices in the \Conditional\ region may participate
only under restricted or conservative rules. Slices in the \Blocked\ region are
excluded from forecasting workflows entirely, preventing structurally hostile
behavior from propagating downstream.

\subsection{Non-Comparative Nature}

FAS classifications are not comparative across slices. A slice classified as
\Conditional\ is not “better” or “worse” than another slice in the same class,
nor does the surface encode relative importance or priority. The admissibility
surface answers a single governance question for each slice in isolation:
\emph{Is forecasting structurally admissible here?}

By framing admissibility as a deterministic surface rather than a score, FAS
avoids conflating structural validity with performance evaluation. This
separation preserves clarity of intent and ensures that admissibility decisions
remain orthogonal to downstream optimization objectives.
